% main.tex — AyShCXR, AIDL-HCSY 2027 (IEEE conference format, max 6 pages)
% Compile on Overleaf: upload the paper/ folder as a zip, set main.tex as the
% main document, press Recompile.
%
% RULE FOR THIS DRAFT: no result number is typed by hand from memory. Every
% \TBD{} is filled only from a results/*.csv file, and the source is noted in a
% comment next to it.
\documentclass[conference]{IEEEtran}
\IEEEoverridecommandlockouts
\usepackage{cite}
\usepackage{amsmath,amssymb}
\usepackage{graphicx}
\usepackage{booktabs}
\usepackage{xcolor}
\usepackage{url}

% visible placeholder: anything red is NOT finished
\newcommand{\TBD}[1]{\textcolor{red}{[#1]}}

% conference rule: no headers, footers or page numbers
\pagestyle{empty}

\begin{document}

\title{Withholding What the Film Cannot Support:\\
Quality-Conditioned Per-Finding Abstention for\\
Multi-Label Chest X-Ray Triage}

% Metadata cannot be changed after final submission: confirm names, order,
% affiliations and whether mentors are co-authors before submitting.
\author{
\IEEEauthorblockN{Subhrakant Sethi}
\IEEEauthorblockA{\textit{Department of Biotechnology} \\
\textit{Thapar Institute of Engineering and Technology}\\
Patiala, India \\
ssethi\_btech24@thapar.edu}
\and
\IEEEauthorblockN{Ayush Singh}
\IEEEauthorblockA{\textit{Department of Biotechnology} \\
\textit{Thapar Institute of Engineering and Technology}\\
Patiala, India \\
asingh3\_btech24@thapar.edu}
}

\maketitle
\thispagestyle{empty}

% ─────────────────────────────────────────────────────────────────────────────
\begin{abstract}
Chest radiographs acquired in primary health centres are often technically
poor: rotated, over- or under-exposed, or blurred. Such faults can make a
deep-learning reader confidently wrong, for example reporting cardiomegaly on a
rotated film. Existing selective-prediction methods for chest X-rays abstain
when the \emph{model} is uncertain, but a degraded film does not necessarily
make the model uncertain. We propose \emph{quality-conditioned per-finding
abstention}: a training-free policy that measures the technical quality of each
film and withholds only those findings that the detected fault is known to
corrupt, leaving the rest of the report intact. On a four-model ensemble
(DenseNet-121, EfficientNet-B4, RAD-DINO, and an NIH-trained DenseNet), under
controlled degradations of NIH ChestX-ray14 test films, uncertainty-based
abstention caught \TBD{x}\% of degradation-induced errors, whereas the
quality-conditioned policy caught \TBD{y}\% at \TBD{equal or better} residual
risk on held-out films. On \TBD{N} real test films, findings that the quality
module marks as compromised \TBD{showed / did not show} a measurable loss of
discrimination. The method needs no retraining and applies to any frozen
classifier.
\end{abstract}

\begin{IEEEkeywords}
chest radiography, selective prediction, abstention, image quality,
uncertainty, deep ensembles, clinical decision support
\end{IEEEkeywords}

% ─────────────────────────────────────────────────────────────────────────────
\section{Introduction}

India has roughly one radiologist per 100{,}000 people, and far fewer in rural
primary health centres (PHCs), where chest radiography is the main imaging
test for tuberculosis, pneumonia, effusion and cardiomegaly. Automated readers
trained on large public datasets \cite{wang2017chestxray8,irvin2019chexpert}
could fill part of this gap, but PHC films are frequently technically poor:
older equipment, briefly trained operators, and patients who cannot hold
position or breath.

A technical fault does not simply add noise; it can \emph{manufacture} or
\emph{hide} specific findings. Rotation widens the projected cardiac
silhouette and mimics cardiomegaly; under-penetration whitens the lung fields
and mimics infiltration; over-penetration and blur erase small nodules. A
classifier has no notion of this and may report the artefact with high
confidence.

The standard safeguard is selective prediction \cite{geifman2017selective}:
abstain when the model is uncertain. For chest X-rays this has been done with
entropy- and confidence-interval-based rejection per pathology
\cite{aperstein2025rejection}. But uncertainty and image quality are different
quantities. When a fault produces a plausible-looking but wrong image, every
ensemble member can agree, and no uncertainty-based rule fires. Recent theory
makes the same point for covariate shift in general \cite{heng2026abstain}.

\textbf{Contributions.}
\begin{itemize}
\item A training-free, \emph{per-finding} abstention policy that conditions on
measured film quality: a fault withholds only the findings it physically
corrupts, so a rotated film still reports pneumothorax.
\item A controlled comparison showing how many degradation-induced errors
uncertainty-based abstention misses on a four-model ensemble, with the policy
chosen on one half of the films and reported on the other.
\item Real-film evidence: on NIH test films, an analysis of whether the faults
the quality module detects in practice degrade the findings it claims they
degrade, controlled for view position.
\end{itemize}

% ─────────────────────────────────────────────────────────────────────────────
\section{Related Work}

\textbf{Selective prediction for chest X-rays.} Aperstein et
al.~\cite{aperstein2025rejection} add entropy- and confidence-interval-based
rejection to a DenseNet-121 classifier with global or class-specific
thresholds. Their rejection signal is model uncertainty only. We use their
setting as the baseline that our policy extends.

\textbf{Quality gating.} Diabetic-retinopathy screening systems routinely
reject a whole image as ``ungradable''. TRUST-CXR \cite{salam2026trustcxr}
combines image-quality metrics with calibrated DenseNet-121 outputs in a
reliability-aware workflow for a single task, tuberculosis. In both cases the
decision is made for the whole image. Ours is made per finding in a
multi-label reader, so a fault that corrupts cardiac size does not suppress an
unrelated pneumothorax.

\textbf{Quality-aware diagnosis.} Che et al.~\cite{che2023qualityaware} use
image-quality labels during \emph{training} to make diagnosis more robust to
degradation. Zhu et al.~\cite{zhu2026lowquality} train report generators to be
quality-agnostic. Both change the model. Our policy abstains instead, needs no
retraining, and applies to any frozen classifier.

\textbf{Chest X-ray technical assessment.} Deep and classical methods estimate
rotation \cite{sun2024rotation} and overall technical adequacy
\cite{meng2022cxrquality} for quality assurance and retake decisions. These
estimators are complementary to our work: any of them could replace our
lightweight quality module as the input to the abstention policy.

% ─────────────────────────────────────────────────────────────────────────────
\section{Method}

\subsection{Classifier ensemble and uncertainty}
Four frozen classifiers are combined: DenseNet-121 \cite{huang2017densenet} and
EfficientNet-B4 \cite{tan2019efficientnet} with dual average/max pooling and
RAD-DINO \cite{perezgarcia2025raddino}, all trained on CheXpert Plus
\cite{chambon2024chexpertplus}, plus a DenseNet-121 trained on NIH
ChestX-ray14 \cite{wang2017chestxray8}. Outputs are merged by canonical finding
name with weights proportional to each model's validation AUC, then calibrated
by per-finding temperature scaling \cite{guo2017calibration}. Uncertainty for
finding $f$ combines the mean Monte-Carlo-dropout spread within models
\cite{gal2016dropout} and the spread of model means across the ensemble
\cite{lakshminarayanan2017ensembles}:
\begin{equation}
u_f = \sqrt{\bar{\sigma}_{\text{within},f}^{2} + \sigma_{\text{between},f}^{2}} .
\end{equation}
A finding is reported positive when its calibrated probability exceeds a
per-finding threshold set at a target specificity on healthy films.
\TBD{state: thresholds tuned on 300 healthy NIH films; the 420 test films that
overlap the tuning sample are excluded from every evaluation.}

\subsection{Film quality assessment}
Three checks are computed from image statistics, without any extra network:
\emph{rotation} (left--right asymmetry of lung-field darkness about an
estimated midline), \emph{exposure} (mean intensity, contrast and clipping),
and \emph{sharpness} (variance of the Laplacian). Thresholds were set at
population percentiles of 400 real VinDr-CXR radiographs: \emph{warn} at about
the 5th/95th percentile and \emph{severe} at about the 1st/99th. A fourth check,
inspiration depth, was disabled because the heuristic located the image border
instead of the diaphragm in a large fraction of films. The rotation check
responds to asymmetry, which large unilateral disease also produces; we return
to this in Section~\ref{sec:limits}.

\subsection{Fault-to-finding map}
Each check $c$ is linked to the findings it can corrupt, $A(c)$, following
standard radiographic reasoning: rotation affects cardiac size, mediastinal
width, effusion, atelectasis and pneumothorax; exposure affects nodule, mass,
infiltration, pneumothorax and emphysema; blur affects nodule, mass and
fibrosis. For a film, finding $f$ receives severity
$s_f = \max\{\text{sev}(c) : f \in A(c)\}$, or none.
\TBD{if A3 is done: compare this map with the per-fault AUC drops measured on
real films.}

\subsection{Abstention policies}
Let $p_f$ be the calibrated probability and $u_f$ the uncertainty.
\begin{itemize}
\item \textbf{Uncertainty-only} (baseline): abstain if $u_f \ge 0.16$, or if
$u_f \ge 0.10$ and $p_f$ is borderline, $0.35 \le p_f \le 0.65$.
\item \textbf{Quality, all calls}: additionally abstain if $s_f$ is severe, or if
$s_f$ is warn and $p_f$ is borderline.
\item \textbf{Quality, positive calls}: apply the quality rule only when the
finding is called positive, i.e.\ withhold a positive the film cannot support
and report it as ``cannot confirm on this film''.
\end{itemize}
No policy has a parameter fitted to the evaluation data: all thresholds are
fixed in the deployed code before the experiments.

% ─────────────────────────────────────────────────────────────────────────────
\section{Experiments}

\subsection{Controlled degradations}
\TBD{N} NIH ChestX-ray14 test films are degraded with four faults at two
severities: over-penetration, under-penetration, Gaussian blur and in-plane
rotation. A decision $(\text{film},f)$ is an \emph{induced error} if $f$ is read
correctly on the clean film and incorrectly on the degraded one. For each
policy we report coverage, residual risk (error rate among reported findings),
and induced errors caught, and compare against uncertainty-ranked abstention at
the \emph{same} coverage. Films are split at random into two halves: the policy
is chosen on half~A by a rule fixed in advance, and all reported numbers come
from half~B. Confidence intervals resample films (1000 bootstrap replicates),
since the 14 findings of one film are not independent.

\subsection{Real films}
The quality module is run on all \TBD{25{,}176} NIH test films. For each check
and finding, discrimination (AUC) is compared between films flagged by that
check and films that pass every check, separately for PA and AP views, because
technically poor films are disproportionately portable AP studies of sicker
patients. Intervals resample patients.

% ─────────────────────────────────────────────────────────────────────────────
\section{Results}

\begin{table}[t]
\caption{Abstention policies on held-out films (half~B, controlled
degradations). ``Unc.\ @ same cov.'' is uncertainty-ranked abstention at the
policy's coverage.}
\label{tab:policies}
\centering
\small
\begin{tabular}{lcccc}
\toprule
Policy & Cov. & Risk & Induced caught & Unc.\ @ same cov. \\
\midrule
None              & 1.00 & \TBD{} & 0 & -- \\
Uncertainty-only  & \TBD{} & \TBD{} & \TBD{} & -- \\
Quality, all      & \TBD{} & \TBD{} & \TBD{} & \TBD{} \\
Quality, positive & \TBD{} & \TBD{} & \TBD{} & \TBD{} \\
Combined, all     & \TBD{} & \TBD{} & \TBD{} & \TBD{} \\
Combined, positive& \TBD{} & \TBD{} & \TBD{} & \TBD{} \\
\bottomrule
\end{tabular}
% source: results/quality_policy_summary.csv (split == "B")
\end{table}

\TBD{Text: which policy the half-A rule selected, its half-B numbers with 95\%
CIs, per-fault breakdown, abstention rate on clean films.}

\TBD{Figure: induced errors caught vs.\ coverage, quality vs.\ uncertainty,
black-and-white safe, saved as EPS/PDF.}

\TBD{Table/figure: real-film AUC drop by check and finding, PA and AP; whether
affected findings drop more than unaffected ones.}

% ─────────────────────────────────────────────────────────────────────────────
\section{Discussion and Limitations}\label{sec:limits}
\begin{itemize}
\item \textbf{Synthetic faults.} The controlled experiment uses simulated
degradations; the real-film analysis is observational and can be confounded
by patient acuity even within a view position.
\item \textbf{Heuristic quality module.} The checks are image statistics, not
anatomical landmark detection. The rotation check responds to any left--right
asymmetry, including large effusions and collapse, so on real films a rotation
flag partly reflects disease.
\item \textbf{Labels.} NIH labels are mined automatically from reports and are
noisy \cite{smit2020chexbert,jain2021visualchexbert}; this limits every
error-based metric in this paper.
\item \textbf{Coverage cost.} Withholding findings shifts work to clinicians;
the abstention rate on clean films is reported for this reason.
\item \textbf{Not a diagnostic device.} The system is decision support and has
not been validated prospectively in a PHC.
\end{itemize}

% ─────────────────────────────────────────────────────────────────────────────
\section{Conclusion}
\TBD{Two or three sentences, written last, stating only what the half-B and
real-film results support.}

\bibliographystyle{IEEEtran}
\bibliography{refs}

\end{document}
